\begin{frame}{경사하강법 대 정규방정식}
  \begin{center}
  \small
  \begin{tabular}{@{}p{0.2\textwidth}p{0.36\textwidth}p{0.3\textwidth}@{}}
    \toprule
     & \kb{정규방정식} & \kb{경사하강법} \\
    \midrule
    반복 & 불필요 --- 정확 & 학습률 튜닝, 스텝 수 필요 \\
    계산 비용 & $(X^TX)^{-1}$ --- 특징 개수 $n$ 의
      \textbf{세제곱}에 비례 & $n$ 이 커져도 매 반복 비용은
      \textbf{선형}으로만 증가 \\
    쓰이는 곳 & 특징이 \textbf{적을} 때 & 특징이 \textbf{많을} 때
      (수천$\sim$수만) \\
    \bottomrule
  \end{tabular}
  \end{center}
  \vskip0.8em
  특징이 적을 땐 정규방정식, 많을 땐 경사하강법을 쓰는 이유다.
\end{frame}
